% =====================================================================
%  IEEE-STYLE CONFERENCE REPORT  (self-contained, from the course template)
%  MLE-AI-201 Machine Learning -- Final Project
%  Build:   tectonic report.tex
% =====================================================================
\documentclass[10pt,twocolumn,a4paper]{article}

% ---- Geometry: IEEE-like margins / two narrow columns ----------------
\usepackage[a4paper,top=0.75in,bottom=1.0in,left=0.68in,right=0.68in,columnsep=0.25in]{geometry}

% ---- Times-like font (standard, ships with every TeX distro) ---------
\usepackage{mathptmx}          % Times text + math
\usepackage[T1]{fontenc}
\usepackage[utf8]{inputenc}    % allows accented / Azerbaijani chars

% ---- Common, dependency-free packages --------------------------------
\usepackage{graphicx}
\usepackage{amsmath,amssymb}
\usepackage{booktabs}
\usepackage{array}
\usepackage{titlesec}
\usepackage{caption}
\usepackage{cite}
\usepackage[hidelinks]{hyperref}
\usepackage{url}
\usepackage{xcolor}

% ---- IEEE-style section headings -------------------------------------
\renewcommand{\thesection}{\Roman{section}}
\renewcommand{\thesubsection}{\Alph{subsection}}
\titleformat{\section}[block]{\centering\normalsize\scshape}{\thesection.}{0.5em}{}
\titleformat{\subsection}[block]{\normalsize\itshape}{\thesubsection.}{0.5em}{}
\titlespacing*{\section}{0pt}{1.4ex plus 1ex minus .2ex}{1ex}
\titlespacing*{\subsection}{0pt}{1.0ex plus 1ex minus .2ex}{0.6ex}

% ---- Abstract / Index-Terms in the IEEE boldface-dash style ----------
\newcommand{\abstractkw}[1]{\noindent\textbf{\textit{#1}}}

\captionsetup{font=small,labelsep=period}
\setlength{\parindent}{1em}
\newcommand{\code}[1]{\texttt{#1}}

% =====================================================================
\begin{document}

% ------------------------- TITLE BLOCK --------------------------------
\twocolumn[
  \begin{@twocolumnfalse}
  \begin{center}
    {\LARGE\bfseries Pricing Azerbaijani Real Estate from Scratch:\\Decision Trees and Pegasos SVMs on bina.az Listings\par}
    \vspace{1.2em}
    {\large
      Aliyeva, Nurana\quad Garibli, Sanan\quad Jamalov, Humbat\quad Nematzada, Farah\quad Shukurov, Elsen\par}
    \vspace{0.5em}
    {\itshape MLE-AI-201 Machine Learning, AI Academy --- Cohort II\par}
    \vspace{1.4em}
  \end{center}
  \begin{quote}
    \small
    \abstractkw{Abstract---}%
    We predict real-estate prices from 100,775 bina.az sale listings with
    two models implemented from scratch in NumPy: a CART-style decision
    tree (Gini, entropy and MSE, vectorised split search, four stopping
    rules) and a soft-margin SVM trained with Pegasos, with random Fourier
    features as a non-linear extension. After deduplicating repeat
    scrapes, parsing text fields and removing target-derived columns
    (\texttt{unit\_price}, \texttt{total\_price}), 63,822 listings remain.
    We frame two tasks: regression on log price and a premium/standard
    tier defined by the training median (210,000~AZN). All hyperparameters
    are chosen on validation data, and the models are benchmarked against
    matched scikit-learn models with 5-fold cross-validation and one test
    evaluation. Our tree matches scikit-learn within one standard deviation
    (CV RMSE 0.338 vs.\ 0.337; CV F1 0.917 vs.\ 0.916) and is the best
    model overall, with test RMSE 0.332 on log price ($R^2=0.80$) and test
    F1 0.929. Our linear SVM matches scikit-learn's hinge-loss solvers (test
    F1 0.843), but a linear decision rule cannot express the interactions
    that drive price. Errors concentrate in the cheapest and most expensive
    listings and near the tier threshold. Code:
    \href{https://github.com/team-M503-SWE-06/ml_final_project}{\texttt{github.com/team-M503-SWE-06/ml\_final\_project}}
    (tag \texttt{v1.0-final}).

    \vspace{0.6em}
    \abstractkw{Index Terms---}decision tree, support-vector machine,
    Pegasos, random Fourier features, regression, classification,
    real-estate price prediction.
  \end{quote}
  \vspace{1.0em}
  \end{@twocolumnfalse}
]

% =====================================================================
\section{Introduction}
Online listings are the main price signal on the Azerbaijani housing
market, but they are noisy: the same flat is re-posted many times, prices
come in two currencies, and numbers arrive as free text. We use one such
dump, the bina.az sale listings, to study two models that we implemented
ourselves in NumPy and benchmark against scikit-learn~\cite{sklearn}.

From one dataset we frame two tasks. \emph{Task A} is regression: predict
the sale price, modelled as $\log(\text{price})$ because prices are heavy
tailed. \emph{Task B} is classification: decide whether a listing is
\emph{premium} (price above the training-set median) or \emph{standard}.
Our contributions are:
\begin{itemize}\setlength\itemsep{0pt}
  \item a leakage-free cleaning pipeline and a frozen, stratified
        70/15/15 split (Section~\ref{sec:data});
  \item a CART-style decision tree~\cite{cart} for both tasks with Gini,
        entropy and MSE criteria, a vectorised split search and four
        stopping rules, tested node by node against scikit-learn;
  \item a Pegasos SVM~\cite{pegasos} with margin and support-vector
        reporting, mini-batches, iterate averaging, and random Fourier
        features~\cite{rff} as the non-linear extension;
  \item a protocol where every choice is made on validation data, a
        5-fold CV benchmark against matched scikit-learn models, one final
        test evaluation, and an error analysis of where the models fail.
\end{itemize}
Everything is reproduced by one command, \code{python -m src.run\_all}.

\section{Background}
\textbf{Decision trees.} A regression or classification tree recursively
splits the data with axis-aligned questions $x_j \le t$. Each split is
chosen to maximise the impurity decrease
\begin{equation}
  \Delta I = I(\text{parent}) - \tfrac{N_L}{N} I(L) - \tfrac{N_R}{N} I(R),
\end{equation}
with Gini $1-\sum_k p_k^2$ or entropy $-\sum_k p_k\log_2 p_k$ for
classification and the variance $\frac{1}{N}\sum_i (y_i-\bar y)^2$ for
regression~\cite{cart}. Deep trees have low bias and high variance, so the
depth or the node sizes must be controlled.

\textbf{Soft-margin SVM and Pegasos.} A linear SVM minimises the
regularised hinge loss
\begin{equation}
  J(\mathbf{w},b)=\tfrac{\lambda}{2}\lVert\mathbf{w}\rVert^{2}
  +\tfrac{1}{n}\sum_{i=1}^{n}\max\!\big(0,\,1-y_i(\mathbf{w}^{\top}\mathbf{x}_i+b)\big),
  \label{eq:svm}
\end{equation}
whose margin width is $2/\lVert\mathbf{w}\rVert$. Pegasos~\cite{pegasos}
solves (\ref{eq:svm}) by stochastic sub-gradient steps with learning rate
$\eta_t = 1/(\lambda t)$, so no quadratic-programming solver is needed.
Random Fourier features~\cite{rff} map $\mathbf{x}$ to
$\phi(\mathbf{x})$ with $\phi(\mathbf{x})^\top\phi(\mathbf{x}')\approx
\exp(-\gamma\lVert\mathbf{x}-\mathbf{x}'\rVert^2)$, so a linear solver on
$\phi$ approximates an RBF-kernel SVM.

% =====================================================================
\section{Data and Preprocessing}\label{sec:data}
We use the Kaggle \emph{Bina.az Sale Project} dump
(\code{sehriyarmemmedli/binaaz-sale-project}, file \code{house\_sale.csv},
100,775 rows $\times$ 51 columns). Every cleaning decision is logged by
\code{data\_prep.clean} and printed by \code{run\_all}:
\begin{itemize}\setlength\itemsep{0pt}
  \item \textbf{Duplicates.} 36,321 rows are repeat scrapes of the same
        listing (same URL); we keep the latest. No exact duplicates remain.
  \item \textbf{Currency and types.} All remaining prices are in AZN (the
        USD path converts at 1.70). Area strings such as ``85~m\textsuperscript{2}''
        and ``6 sot'' (1 sot = 100~m\textsuperscript{2}) and floors such as
        ``3~/~9'' are parsed into numbers (\code{floor}, \code{total\_floors}).
  \item \textbf{Outliers.} We drop prices below 1,000~AZN or outside a
        fixed 10,000--5,000,000~AZN range (457 rows), impossible areas and
        per-category area outliers (175 rows), and set coordinates outside
        Azerbaijan's bounding box to missing. Limits are fixed domain
        values, not data quantiles, so they cannot leak test information.
  \item \textbf{Leakage.} \code{unit\_price} ($=$ price / area) and
        \code{total\_price} ($=$ price in 99.99\% of rows) are arithmetic
        functions of the target and are removed, together with identifiers
        (\code{owner\_name}, \code{shop\_name}, \code{address}, URLs, ids),
        free text and scrape metadata (views, timestamps).
  \item \textbf{Missing values.} Medians are computed on the training rows
        only and every imputed column gets a 0/1 \code{\_missing} flag
        (\code{land\_area} is 85\% missing, mostly flats).
\end{itemize}
The clean table has 63,822 listings. The feature matrix has 72 columns:
7 numeric (area, rooms, floor, total floors, land area, latitude,
longitude), 6 missing-value flags, 6 binary flags (repair, bill of sale,
mortgage, agent, residential complex, Baku) and 53 one-hot levels for
category, district and location; rare levels learned on the training rows
are pooled into ``other''.

\begin{table}[t]
  \centering
  \caption{Dataset summary after cleaning.}
  \label{tab:data}
  \begin{tabular}{@{}lr@{}}
    \toprule
    Property & Value \\
    \midrule
    Raw rows / columns            & 100,775 / 51 \\
    Rows after cleaning           & 63,822 \\
    Features used                 & 72 \\
    Train / Val / Test            & 44,676 / 9,573 / 9,573 \\
    Tier threshold (train median) & 210,000 AZN \\
    Premium share (each split)    & 49.3\% \\
    \bottomrule
  \end{tabular}
\end{table}

The split is a seeded, stratified 70/15/15 partition frozen to
\code{data/split\_indices.npz} (Table~\ref{tab:data}). The tier label is
\emph{premium} when the price exceeds the training-set median of
210,000~AZN; validation and test labels reuse that threshold, so the test
set never influences it. The resulting classes are balanced (49.3\%
premium in every split).

% =====================================================================
\section{Method}
\subsection{Decision tree (from scratch)}
\code{src/decision\_tree.py} implements one \code{DecisionTree} class for
both tasks (Gini, entropy or MSE).

\textbf{Split search.} At each node every feature is sorted once.
A threshold between sorted positions $i$ and $i+1$ sends the first $i+1$
rows left, so cumulative sums over the sorted rows give the class counts
(classification) or $\sum y$ and $\sum y^2$ (regression) of every left
child at once; the right child is ``node total minus left''. This scores
all candidate thresholds of all features in one vectorised pass, costing
$O(nd\log n)$ per node instead of $O(n^2 d)$. Candidates are midpoints
between consecutive \emph{distinct} values, as in scikit-learn.

\textbf{Growth and stopping.} The tree grows by recursive partitioning
(driven by an explicit stack, so very deep trees cannot hit Python's
recursion limit). A node becomes a leaf when it reaches
\code{max\_depth}, has fewer than \code{min\_samples\_split} rows, is pure,
has no threshold leaving \code{min\_samples\_leaf} rows on each side, or
when the weighted decrease $(N_t/N)\,\Delta I$ is below
\code{min\_impurity\_decrease} (scikit-learn's definition). Leaves store
class frequencies (majority vote and \code{predict\_proba}) or the mean.
Feature importances are the normalised total weighted impurity decrease.

\textbf{Ties and validation.} When two features give the same gain we take
the lowest feature index; scikit-learn breaks such ties at random. The
tree rejects NaN or infinite inputs, checks the column count at predict
time, and computes thresholds as $a/2+b/2$ to avoid overflow.

\textbf{Tests.} \code{tests/test\_decision\_tree.py} compares our tree node
by node (feature, threshold, samples, value) with scikit-learn on a toy
set and on seeded random data for all three criteria; each example is
first checked to give the same scikit-learn tree for ten random seeds, so
the expected tree is fully determined. Further tests check that each
stopping rule changes the tree, that \code{predict\_proba} sums to one, and
the input validation.

\subsection{Support-vector machine (from scratch)}
\code{src/svm.py} implements \code{PegasosSVM}. Labels are mapped to
$\pm1$; the bias is treated as the weight of a constant-1 feature, so the
objective actually minimised is (\ref{eq:svm}) with
$\tfrac{\lambda}{2}(\lVert\mathbf{w}\rVert^2+b^2)$, a negligible change
on standardised features that keeps the very large first steps
($\eta_1 = 1/\lambda$) stable. Each step over a mini-batch $A_t$ of size
$k$ with active set $A_t^+=\{i: y_i(\mathbf{w}^\top\mathbf{x}_i+b)<1\}$ is
\begin{equation}
  \mathbf{w} \leftarrow (1-\eta_t\lambda)\,\mathbf{w}
  + \tfrac{\eta_t}{k}\textstyle\sum_{i\in A_t^+} y_i\mathbf{x}_i ,
\end{equation}
with three schedules: Pegasos $1/(\lambda t)$, $\eta_0/\sqrt{t}$ and
constant. An optional projection onto $\lVert(\mathbf{w},b)\rVert\le
1/\sqrt{\lambda}$ follows each step, and the returned model averages the
iterates of the final epoch (tail averaging), which removes most of the
step noise. The objective is logged every epoch, and the model reports
$\lVert\mathbf{w}\rVert$, the margin width $2/\lVert\mathbf{w}\rVert$ and
the support vectors ($y_i f(\mathbf{x}_i)\le1$). Features are
standardised with training statistics first, because the hinge loss is
scale sensitive.

\textbf{Non-linear extension.} We chose random Fourier features
(\code{RFFPegasosSVM}): $\phi(\mathbf{x})=\sqrt{2/D}\cos(W^\top\mathbf{x}+\mathbf{c})$
with $W_{jk}\sim\mathcal{N}(0,2\gamma)$, $c_k\sim U(0,2\pi)$ and $D=500$,
followed by the unchanged linear solver. Kernelised Pegasos would need
$O(n^2)$ kernel evaluations per epoch (about $2\times10^9$ for our 44,676
training rows) and $O(n)$ work per prediction; RFF costs $O(nD)$ per epoch
and reuses the tested linear code.

% =====================================================================
\section{Experimental Setup}
\textbf{Metrics.} Task A: RMSE, MAE and $R^2$ on $\log(\text{price})$,
plus RMSE and MAE in AZN. Because $\exp(\widehat{\log y})$ estimates the
median rather than the mean, AZN predictions are rescaled by Duan's
smearing factor~\cite{duan}, estimated on validation residuals. Task B:
precision, recall, F1 (premium is the positive class), ROC-AUC and the
confusion matrix.

\textbf{Protocol.} Hyperparameters are chosen on the validation split
only: depth, criterion and leaf-size sweeps for the tree, a 5-fold CV
search on the training split for the final tree settings, and $\lambda$,
epoch and schedule sweeps (and the RFF bandwidth $\gamma$) for the SVM.
The benchmark then reports 5-fold CV mean $\pm$ std on the training split
and evaluates every model once on the test split. Every random step uses
seed 42 and the split is frozen, so reruns give identical numbers.

\textbf{Baselines.} scikit-learn models with matched settings:
\code{DecisionTreeRegressor/Classifier} with the same hyperparameters;
for the SVM, \code{SGDClassifier} (hinge, $\alpha=\lambda$) and
\code{LinearSVC} (hinge, $C=1/(\lambda n)$), both with the scaler inside a
\code{Pipeline}, and an RBF \code{SVC} on a 10,000-row training subsample
(SVC is $O(n^2)$).

% =====================================================================
\section{Results}
\subsection{Tuning the tree}
Fig.~\ref{fig:depth} shows the bias--variance trade-off of the regression
tree: training RMSE keeps falling with depth, while validation RMSE is
lowest at depth~10 (0.346) and rises again (0.366 at depth~20, where the
training RMSE is 0.095). For the tier task, Gini and entropy are within
0.003 F1 at every depth; Gini peaks at depth~15 (validation F1 0.919), and
larger \code{min\_samples\_leaf} values only lowered validation F1.

A single depth limit cuts every branch at the same level, although dense
districts need deeper trees than sparse ones. A 5-fold CV search on the
training split over all four stopping rules therefore replaced it with
node-size rules: \code{min\_samples\_split}=75 with no depth limit for
Task~A (CV RMSE $0.3621\pm0.0033 \rightarrow 0.3377\pm0.0032$) and Gini
with \code{max\_depth}=20 and \code{min\_impurity\_decrease}=$2\cdot10^{-5}$
for Task~B (CV F1 $0.9115\pm0.0027 \rightarrow 0.9166\pm0.0019$). Both
gains hold on every fold, and these are the configurations used below.

\begin{figure}[t]
  \centering
  \includegraphics[width=0.92\linewidth]{figures/tree_reg_depth.png}
  \caption{Regression tree, depth sweep (validation): the gap between
  training and validation RMSE opens after depth~10.}
  \label{fig:depth}
\end{figure}

\subsection{Tuning the SVM}
On standardised features, validation F1 of the linear Pegasos SVM is flat
for $\lambda\in[10^{-6},10^{-3}]$ (0.842--0.844) and drops as the
regularisation grows (0.791 at $\lambda=1$), while the margin
$2/\lVert\mathbf{w}\rVert$ widens from 0.007 to 3.87 and the number of
support vectors rises from 7,081 to 36,837 (Fig.~\ref{fig:lambda}). We
keep $\lambda=10^{-3}$ (F1 0.844, margin 0.71). The objective decreases
over the epochs; the Pegasos and $\eta_0/\sqrt t$ schedules converge to
the same objective (about 0.385 after 50 epochs), while a constant step of
0.01 oscillates and stays near 0.45. Averaging the last epoch reaches, in
10 epochs, the objective that the last iterate needs about 50 epochs for
(0.385 vs.\ 0.395 at 10 epochs). Our weight vector points in almost the same
direction as scikit-learn's solutions (cosine 0.997 with both
\code{SGDClassifier} and \code{LinearSVC}). Random Fourier features did
not help: the best bandwidth ($\gamma=0.01$, $D=500$) reached validation
F1 0.809, below the linear model.

\begin{figure}[t]
  \centering
  \includegraphics[width=0.92\linewidth]{figures/svm_lambda_sweep.png}
  \caption{Pegasos SVM, $\lambda$ sweep (validation): stronger
  regularisation widens the margin but costs accuracy beyond
  $\lambda=10^{-3}$.}
  \label{fig:lambda}
\end{figure}

\begin{figure*}[t]
  \centering
  \includegraphics[width=\textwidth]{figures/error_analysis.png}
  \caption{Test-set error analysis of the final regression tree:
  predictions vs.\ actual log price, RMSE by price quintile, and RMSE by
  district frequency (the sparse group has only nine listings).}
  \label{fig:errors}
\end{figure*}

\subsection{Benchmark against scikit-learn}
Table~\ref{tab:cv} gives 5-fold CV on the training split and
Table~\ref{tab:results} the single test evaluation. Our tree matches
scikit-learn's within one CV standard deviation on both tasks. Our SVM is
on par with the matched linear baselines (CV F1 0.843 vs.\ 0.840 and
0.842); only the RBF \code{SVC}, trained on a 10,000-row subsample, is
slightly higher (0.850 $\pm$ 0.009). On the test set the tree reaches an
RMSE of 0.332 on log price ($R^2=0.80$) and F1 0.929 on the tier task,
against F1 0.843 for the SVM.

The trees are structurally close to scikit-learn's: for Task~B both have
depth~20 and 1,655 vs.\ 1,661 leaves, and 99.1\% of test predictions are
identical; for Task~A (no depth limit) the depths are 38 vs.\ 33 with
1,507 vs.\ 1,497 leaves and 94.2\% identical predictions. The gaps come
from tie-breaking (we take the lowest feature index; scikit-learn draws
features at random) and from scikit-learn storing thresholds in float32;
on tie-free examples our unit tests reproduce its trees node by node.
Speed is the main difference: our tree fits in 6.9~s (Task~A) and 17.5~s
(Task~B) per fold against 0.7--0.8~s, and our SVM in 8.8~s against 0.6~s
for \code{SGDClassifier}, because the split search and the Pegasos steps
run in Python/NumPy instead of compiled code.

\begin{table}[t]
  \centering
  \caption{5-fold CV on the training split (mean $\pm$ std) and fit time per fold.}
  \label{tab:cv}
  \small\setlength{\tabcolsep}{4pt}
  \begin{tabular}{@{}llcr@{}}
    \toprule
    Task & Model & Score & Fit (s) \\
    \midrule
    A (RMSE) & Ours: tree            & $0.3377\pm0.0032$ & 6.87 \\
             & sklearn tree           & $0.3369\pm0.0038$ & 0.74 \\
    \midrule
    B (F1)   & Ours: tree            & $0.9166\pm0.0019$ & 17.50 \\
             & sklearn tree           & $0.9162\pm0.0030$ & 0.77 \\
             & Ours: Pegasos SVM     & $0.8433\pm0.0029$ & 8.76 \\
             & sklearn SGD (hinge)    & $0.8404\pm0.0036$ & 0.60 \\
             & sklearn LinearSVC      & $0.8416\pm0.0022$ & 1.83 \\
             & sklearn RBF SVC$^\dagger$ & $0.8497\pm0.0091$ & 5.43 \\
    \bottomrule
    \multicolumn{4}{@{}l}{\footnotesize $^\dagger$ 10,000-row subsample per fold.}
  \end{tabular}
\end{table}

\begin{table}[t]
  \centering
  \caption{Headline results on the test set (one evaluation).}
  \label{tab:results}
  \begin{tabular}{@{}lcccc@{}}
    \toprule
    \multicolumn{5}{@{}l}{\textit{Task A: log(price)}} \\
    Model & RMSE & MAE & $R^2$ & \\
    \midrule
    Ours: tree   & 0.332 & 0.193 & 0.802 & \\
    sklearn tree & 0.332 & 0.193 & 0.802 & \\
    \midrule
    \multicolumn{5}{@{}l}{\textit{Task B: price tier (premium = positive)}} \\
    Model & Prec. & Recall & F1 & AUC \\
    \midrule
    Ours: tree        & 0.926 & 0.932 & 0.929 & 0.925 \\
    sklearn tree      & 0.927 & 0.933 & 0.930 & 0.927 \\
    Ours: SVM         & 0.823 & 0.865 & 0.843 & 0.922 \\
    sklearn SGD       & 0.828 & 0.849 & 0.839 & 0.921 \\
    sklearn LinearSVC & 0.828 & 0.856 & 0.842 & 0.921 \\
    \bottomrule
  \end{tabular}
\end{table}

\subsection{Feature importance and SVM weights}
The tree's impurity-based importances are dominated by floor area,
followed by location (longitude, latitude and the pooled ``other''
location), land area and the category (land, house, new build). The SVM
weights (standardised features) tell a consistent story: the largest
positive weights are rooms, area, the missing-rooms flag (typical of land
plots and houses), \emph{new build}, and central districts and
neighbourhoods (Sabail, White City, Shah Ismail Khatai); the most negative
are peripheral districts (Khachmaz, Absheron, Surakhani), \emph{land} and
\emph{old build}.

\subsection{Error analysis}
Fig.~\ref{fig:errors} and Table~\ref{tab:bands} break the test errors of
the final trees down. Task~A errors are U-shaped across price quintiles:
the middle three quintiles have an RMSE of 0.22--0.25 (median error
9--11\%), but the cheapest quintile reaches 0.40 and the most expensive
0.47. In the top 5\% of prices the RMSE is 0.75 and the median error is
44\%. The ten largest errors are mostly land plots whose prices range from
20,000 to 5,000,000~AZN for similar areas and districts, which the
features cannot separate. Task~B errors concentrate at the tier boundary:
listings within 10\% of the 210,000~AZN threshold are misclassified
25.4\% of the time, against 4.8\% for the rest, and the middle price
quintile has a 20.3\% error rate against 1.7--6.2\% elsewhere. Only nine
test listings fall in sparse or missing districts; their RMSE is high
(1.50) but too few to draw conclusions from.



\begin{table}[t]
  \centering
  \caption{Test errors of the final trees by price quintile.}
  \label{tab:bands}
  \begin{tabular}{@{}lrccc@{}}
    \toprule
    & & \multicolumn{2}{c}{Task A} & Task B \\
    Quintile & $n$ & RMSE & med.\ err.\ (\%) & error rate \\
    \midrule
    Q1 cheapest & 1,851 & 0.398 & 13.5 & 0.029 \\
    Q2          & 1,993 & 0.221 &  8.9 & 0.045 \\
    Q3          & 1,862 & 0.228 &  8.7 & 0.203 \\
    Q4          & 1,858 & 0.250 & 10.7 & 0.062 \\
    Q5 luxury   & 2,009 & 0.474 & 19.2 & 0.017 \\
    \bottomrule
  \end{tabular}
\end{table}

% =====================================================================
\section{Discussion and Limitations}
\textbf{Why the tree beats the SVM.} Price depends on interactions such as
area $\times$ district $\times$ category, which a tree models with nested
splits. The linear SVM can only add up per-feature contributions, so it
plateaus at F1~0.84 regardless of $\lambda$, and its errors do not shrink
with more epochs. Random Fourier features should capture interactions,
but with $D=500$ components over 72 mostly binary features the
approximation was too coarse and underperformed the linear model; a
kernel SVC on 10,000 rows gained less than one point. A tree, in
contrast, overfits easily, which is why the node-size rules mattered.

\textbf{Where the models fail.} Both extremes of the price range are hard.
A leaf predicts the mean log price of similar listings, so predictions are
pulled towards the middle: cheap listings are over-estimated and luxury
ones under-estimated. On top of this, $\exp(\widehat{\log y})$ estimates
the median, not the mean (Jensen's inequality); the smearing factor
(1.059) corrects the average level in AZN but cannot fix the tails. The
largest single errors are land plots, whose price depends on things the
data does not describe well (zoning, exact location, road access). For
the tier task, most mistakes are listings priced close to the threshold,
which is expected for a hard 210,000~AZN cut.

\textbf{Limitations.} (i) Errors in AZN remain large (test RMSE
253,669~AZN, MAE 79,550~AZN) because a few multi-million listings dominate
the squared error; the log scale is the fairer summary. (ii) The data is a
single scrape of asking prices, not transaction prices, and repeated
listings had to be deduplicated by URL. (iii) One-hot encoding lets each
split separate only one district from the rest; latitude and longitude
partly compensate. (iv) The 5-fold CV search that chose the final tree
settings is not part of \code{run\_all}; the benchmark re-runs CV for the
chosen settings only. (v) Our implementations are 9--23$\times$ slower to
fit than scikit-learn's compiled ones. (vi) The SVM regularises the bias
as a constant feature, a small departure from the textbook objective.

\section{Conclusion}
A from-scratch NumPy decision tree and Pegasos SVM, trained on a cleaned,
leakage-free version of the bina.az listings, match scikit-learn within
one CV standard deviation. The tree is the better model for both tasks:
test RMSE 0.332 on log price ($R^2=0.80$) and F1 0.929 on the price tier,
against F1 0.843 for the linear SVM. The remaining errors sit at the price
extremes and at the tier boundary. Next steps are ensembles of our own
trees to reduce the variance of a single tree, target or ordered encoding
for districts, and richer location features.

\section*{Tools and Acknowledgements}
We used the Kaggle \emph{Bina.az Sale Project} dataset, NumPy, pandas,
Matplotlib, and scikit-learn for the baselines only. AI assistance: Claude
(Anthropic) was used to test the initial versions drafted according to each member's role. All code was
reviewed in pull requests and checked against scikit-learn.


% =====================================================================
\begin{thebibliography}{9}
\bibitem{cart}
L.~Breiman, J.~Friedman, R.~Olshen, and C.~Stone,
\emph{Classification and Regression Trees}. Wadsworth, 1984.
\bibitem{pegasos}
S.~Shalev-Shwartz, Y.~Singer, N.~Srebro, and A.~Cotter,
``Pegasos: primal estimated sub-gradient solver for SVM,''
\emph{Mathematical Programming}, vol.~127, no.~1, pp.~3--30, 2011.
\bibitem{rff}
A.~Rahimi and B.~Recht, ``Random features for large-scale kernel
machines,'' in \emph{Advances in Neural Information Processing Systems},
2007.
\bibitem{duan}
N.~Duan, ``Smearing estimate: a nonparametric retransformation method,''
\emph{Journal of the American Statistical Association}, vol.~78,
no.~383, pp.~605--610, 1983.
\bibitem{sklearn}
F.~Pedregosa \emph{et al.}, ``Scikit-learn: machine learning in Python,''
\emph{Journal of Machine Learning Research}, vol.~12, pp.~2825--2830, 2011.
\end{thebibliography}

\end{document}
